% ==============================================================================
% CONCLUSIONES INDIVIDUALES
% ==============================================================================
% Instrucciones de la plantilla:
% Formato: #CUENTA: CONCLUSIÓN
% Cada integrante del equipo debe redactar de forma personal, concreta y 
% coherente sus conclusiones individuales sobre la práctica y los conceptos 
% gráficos aprendidos.
% ==============================================================================

\section{Conclusiones Individuales}

\subsection*{\#321137135: Basilio Illescas Andrés}
% Redactar la conclusión individual del alumno aquí:
A través de las actividades implementadas, se logró afianzar el conocimiento sobre...

\subsection*{\#424065465: Cruz Macedo Samuel Santiago}
% Redactar la conclusión individual del alumno aquí:
El análisis de los experimentos realizados permitió identificar...

\subsection*{\#321258173: Medel Piña Ángel de Jesús}
% Redactar la conclusión individual del alumno aquí:
La práctica permitió contrastar la teoría vista en clase con la implementación en OpenGL...

\subsection*{\#321150503: Pérez Olvera Alexis Abraham}
% Redactar la conclusión individual del alumno aquí:
Durante el desarrollo de la práctica se comprendieron los conceptos fundamentales de...

\subsection*{\#320224159: Segura Cedeño Luisa María}
En esta práctica logramos entender cómo funcionan los modelos de iluminación, pasando de ver objetos planos a darles una apariencia mucho más real. Lo que más me llamó la atención fue ver cómo se combinan las tres componentes principales: la ambiental para la luz general, la difusa para las zonas que reciben luz directa y la especular para el brillo que pueden simular materiales como vidrio, espejo, agua, entre otros. Además, aprendí que configurar bien los vectores normales y las propiedades de la fuente de luz es importante para que los objetos no se vean opacos ni planos. Por otro lado, añadimos audio a la escena; aunque es un procedimiento sencillo de implementar, hace mucha diferencia en la experiencia y hace que el ambiente se sienta todavía más realista e inmersivo.
